% 第 19 题四边形图：按固定源内容流矢量层精确值复刻（1 pt → 0.03515 cm，与全库同制）。
% 原点取 A 的 x=272.6508 pt、底缘（A/B 标签下缘）y（顶向下）=630.2490 pt；y 向上（Y = 630.2490 − y_页）。
% 源实测（figure-inventory §4）：4 条边＝4 条独立路径（全实线、无填充、无箭头、线宽 0.3985；对角线不画）。
%   顶点 A(272.6508,619.6931)、B(391.7075,619.6931)、C(388.5344,560.1628)、D(315.1121,530.0029)；
%   AB 严格水平（两端 y 逐位相同）；BC 接近竖直（Δx=−3.1731）。
%   实长 DA=99.2335 / DC=79.3754 / BA=119.0567 / BC=59.6148，折 cm 为 3.5007/2.8002/4.2000/2.1031，
%   以 0.7 cm 为单位即 5.0010/4.0003/6.0001/3.0044 ≈ 5/4/6/3，与题干 AB=6、BC=3、CD=4、AD=5
%   逐一对应（偏差 ≤0.0015 单位）⇒ 本图按 0.7 cm/单位按比例作图、角度亦按真值，重绘须保留（S-3）。
\begin{tikzpicture}[x=0.03515cm,y=0.03515cm,every node/.style={font=\small,inner sep=0pt},
  gl/.style={line width=0.3985bp,line join=round}]
% ===== 四条边（各为独立路径，照源写法）=====
\draw[gl] (42.4613,100.2461)--(0,10.5559);          % DA
\draw[gl] (42.4613,100.2461)--(115.8836,70.0862);   % DC
\draw[gl] (119.0567,10.5559)--(0,10.5559);          % BA
\draw[gl] (119.0567,10.5559)--(115.8836,70.0862);   % BC
% ===== 标签 4 个（照源字形盒中心定位）=====
\node at (-4.9333,4.6620) {\(A\)};
\node at (123.7242,4.6620) {\(B\)};
\node at (120.4372,74.0540) {\(C\)};
\node at (42.3167,108.2510) {\(D\)};
\end{tikzpicture}
